% Generated draft: requires mathematical and source review.
\documentclass[11pt]{article}
\usepackage[margin=25mm]{geometry}
\usepackage{fontspec}
\setmainfont{Latin Modern Roman}
\usepackage{amsmath,amssymb,mathtools,unicode-math}
\setmathfont{Latin Modern Math}
\usepackage{enumitem,needspace,xurl,hyperref}
\setlength{\parindent}{0pt}
\setlength{\parskip}{5pt}
\setlength{\emergencystretch}{3em}
\newcommand{\N}{\mathbb{N}}
\newcommand{\Z}{\mathbb{Z}}
\newcommand{\Q}{\mathbb{Q}}
\newcommand{\R}{\mathbb{R}}
\newcommand{\C}{\mathbb{C}}
\newcommand{\F}{\mathbb{F}}
\newcommand{\A}{\mathbb{A}}
\newcommand{\PP}{\mathbb{P}}
\newcommand{\E}{\mathbb{E}}
\newcommand{\eps}{\varepsilon}
\newcommand{\dd}{\,\mathrm{d}}
\begin{document}
\section*{UnsolvedMath ID: 30001082}
\textbf{Source problem code:} catalogue entry
\subsection{Definitions and mathematical statement}
A Cox process $N$ directed by a random locally finite measure $\Lambda$ is a point process which, conditional on $\Lambda$, is Poisson with intensity measure $\Lambda$. Thus counts on disjoint measurable sets are conditionally independent and $N(B)\mid\Lambda$ has Poisson distribution with mean $\Lambda(B)$.

In a conventional log-Lévy-driven construction on a Euclidean observation space, a Lévy random basis $L$ is an independently scattered infinitely divisible random measure, a field is formed as
\[
Y_\theta(x)=\int k_\theta(x,s)\,L_\theta(ds),\qquad
\Lambda_\theta(B)=\int_B e^{Y_\theta(x)}\,dx,
\]
and $N$ is the Cox process directed by $\Lambda_\theta$. This is explanatory notation, not a specification recovered from the source: existence of the stochastic integral and local finiteness of $\Lambda_\theta$ require appropriate assumptions, which are not supplied.

The source task is to develop statistical inference for log-Lévy-driven Cox processes, including their driving intensity and parameters. It does not choose an estimation target or observation regime. ``Driving intensity'' may mean the latent random intensity, its mean function, or characteristics of the underlying Lévy driver; these are distinct inferential targets. Likewise $\theta$ may encode the kernel and driver law, but no particular parameterization, finite-dimensionality or identifiability statement is given. Valid methods and their guarantees therefore require a specified model, observed data, target, regularity conditions and sampling regime. No estimator, likelihood approximation, rate or asymptotic distribution is imposed by the record.

The supplied triage checked 2026-08-21 labels this an open research direction. It says the model class and moment and mixing properties have been developed but general inference remains open for Lévy-driven and log-Lévy-driven Cox processes. The broader assessment is retained without conflating direct Lévy-driven intensities with the log-Lévy-driven construction asked for in the main statement. The registry also records manual review labeling the extracted research request exact and coherent.
\subsection{Short English statement}
Develop inference for a Cox point process with log-Lévy-driven intensity, covering its intensity and model parameters. The driver, kernel, observations and distinction between latent-intensity recovery and parameter estimation must be specified before a precise statistical theorem can be stated.
\subsection{Sources}
\begin{itemize}
\item Local UnsolvedMath registry, permanent ID 30001082, source code catalogue entry; supplied complete record and exact-statement manual-review note. Literature checked 2026-08-21: the assessment is supplied dated curation, not independently verified here.
\item New Perspectives in Stochastic Geometry (2009), Oberwolfach Reports, DOI 10.4171/owr/2008/47, pp. 2655--2702: \url{https://doi.org/10.4171/owr/2008/47}. Proposed-year metadata gives 2008.
\item Lévy-based Cox point processes: \url{https://doi.org/10.1239/aap/1222868178}.
\end{itemize}
\end{document}
